% !TEX root = CFAST_Validation_Guide.tex

\chapter{Pressure}

Comparisons between measurement and prediction of compartment pressure for the NIST/NRC test series, two of the NBS furniture tests, LLNL Enclosure, and PRISME are shown in Appendix~\ref{sec:Graphs}.  Figure \ref{fig:Pressure_Scatter} shows a comparison of predicted and measured values for compartment pressure.
\label{Compartment Over-Pressure}

\begin{figure}
\begin{center}
\includegraphics[height=4.0in]{SCRIPT_FIGURES/Scatterplots/Compartment_Pressure}  \\
\end{center}
\caption[Summary, Compartment Pressure]{Summary, Compartment Pressure.} \label{fig:Pressure_Scatter}
\end{figure}

For those tests in which the door to the compartment is open, the magnitude of the pressures are only a few Pascals; however, when the door is closed, the over-pressures are several hundred Pascals. In general, prediction of pressure in CFAST in closed compartments is critically dependent on correct specification of the leakage from the compartment.  Compartments are rarely entirely sealed, and small changes in the leakage area can produce significant changes in the predicted over-pressure. Information on leakage is not typically available for experimental data sets. For closed compartment tests, which typically included mechanical ventilation systems, leakage was estimated by matching the experimental pressure prior to ignition by inserting a small vent in the wall of each compartment and comparing the measured and calculated compartment pressure without a fire in the simulation.

